
\documentclass[a4paper,12pt]{article}

\input{glyphtounicode}
\pdfgentounicode=1

\usepackage{mathptmx}
\usepackage{fancyhdr}
\usepackage{hyperref}
\usepackage[english]{babel}
\usepackage[utf8]{inputenc}
\usepackage[version=4]{mhchem}
\usepackage[usenames,dvipsnames]{color}
\usepackage[
  a4paper,
  top    = 2.00cm,
  bottom = 2.00cm,
  left   = 1.70cm,
  right  = 1.70cm
]{geometry}

\renewcommand{\baselinestretch}{1.5}
\renewcommand{\headrulewidth}{0pt}
\renewcommand{\footrulewidth}{0pt}

\hypersetup{
    colorlinks=true,
    linkcolor=RoyalBlue,
    urlcolor=RoyalBlue,
    citecolor=RoyalBlue,
    anchorcolor=RoyalBlue
}

\urlstyle{same}
\pagestyle{fancy}
\fancyhf{}
% \fancyfoot[R]{\thepage}

\begin{document}

\begin{center}
    {\Huge \scshape {\fontsize{25}{30}\selectfont{Reza}} {\fontsize{25}{30}\selectfont{\textbf{Aliasgari Renani}}}} \\ \vspace{3pt}
    {\small \raisebox{-0.2\height}\ {06/12/2025}} $|$
    {\small \raisebox{-0.2\height}\ {Letter of Motivation}} $|$
    {\small \raisebox{-0.2\height}\ {Physics}} $|$
    {\small \raisebox{-0.2\height}\ {Marietta Blau Institute of Particle Physics}} $|$
    {\small \raisebox{-0.2\height}\ {PhD Position}}
    \vspace{-10pt}
\end{center}
\vspace{-10pt}
\noindent\rule{\textwidth}{0.5pt}

My research objective is to contribute to the advancement of experimental particle physics through the physics analysis of collision data, in particular by applying knowledge of radiation effects on detectors and creating FPGA-based systems for efficient data processing. The PhD position for Physics Analysis with the CMS experiment at the Marietta Blau Institute of Particle Physics of the Austrian Academy of Sciences offers an excellent opportunity for me to pursue this goal. I am now completing the final year of my Masters degree in Plasma Physics at Moscow Institute of Physics and Technology, where my research on semiconductor devices exposed to different kinds of radiation and FPGA development has strengthened my interest in studying high energy particles. I am especially motivated by the possibility of contributing to analyses within a large international collaboration and using hardware expertise for experimental physics.

\vspace{10pt}
I joined Dr.\ Koveshnikov's laboratory in the Russian Academy of Sciences over two years ago to study radiation effects on microelectronic structures, focusing on electrical characterization of semiconductors and \ce{SiO2}-based MOS devices. We studied the impact of electron and gamma irradiation and hydrogen plasma treatment on MOS devices, identified traps by location (oxide, semiconductor, interface) and carrier type (electrons, holes). This work has resulted in a first-author publication\footnote{R. Aliasgari Renani, O.A. Soltanovich, M.A. Knyazev, S.V. Koveshnikov, Investigation of low energy electron irradiated \ce{SiO2} based MOS devices by C-V and TSC techniques, Russian Microelectronics, 2023} and provides insight into device stability under bias stress, radiation immunity, and oxide trap density at the dielectric-semiconductor interface. I also developed optimized MATLAB applications to automate the experimental characterization techniques, which reduced run-times and created measurement pipelines. 

\vspace{10pt}
After my undergraduate degree, I joined the System-on-Chip Development Laboratory and the Laboratory of Plasma Systems under the supervision of Dr.\ Vasiliev to develop FPGA-based systems and to study radiation and plasma effects on these structures. I implemented video-processing algorithms by manual Verilog coding and by generating HDL from Simulink models. I verified functionality of the algorithms using testbenches and Python automation scripts. In the Plasma Lab, we ran experiments where an FPGA with an attached camera and a synthesized image-processing design was placed in an electron-beam plasma system with low-pressure oxygen and exposed to electron irradiation and X-rays generated from a tungsten target. The FPGA output signal was monitored while thermal cycling was applied and preliminary functionality tests were performed. 

\newpage

\vspace{10pt}
My experience with irradiated semiconductors, plasma systems and FPGA design motivated me to continue my PhD studies at the Marietta Blau Institute of Particle Physics of the Austrian Academy of Sciences. I aim to engage in physics analyses of CMS data to search for signs of new physics or to perform precision measurements of standard model parameters. My background will enable me to contribute to the institutes projects on detector development and upgrades, such as silicon trackers and calorimeters that must withstand high radiation environments, by applying my knowledge of radiation effects and FPGA implementation for data acquisition. I am enthusiastic to work under the guidance of Dr.\ Robert Schöfbeck at the Marietta Blau Institute to integrate analysis techniques with hardware solutions for advancing CMS research.

\vspace{10pt}
Having been an international student and having collaborated across multiple scientific institutions, I am comfortable working between partner laboratories and I would welcome the opportunity to split time between Vienna and CERN to advance this research. I am driven by research that couples numerical methods with experiments, and I look forward to leveraging my skills in FPGA development for real-time data processing and to validate the data against results from the LHC runs, which would ultimately contribute to breakthroughs in understanding fundamental particles. It would be an honor to study and work at the Austrian Academy of Sciences in Austria, a country with longstanding traditions of scientific and cultural excellence.

\end{document}